\providecommand{\openpoint}[1]{\par\noindent\fbox{\parbox{0.97\linewidth}{%
  \textbf{Open point:} #1}}\par}

\chapter{Introduction}
\label{ch:introduction}

\section{Context and motivation}
\label{sec:intro-context}

Compact pressurised water reactors (PWRs) have propelled naval vessels
since the 1950s, and some 700 reactors have been used at sea since
then~\cite{wna_ships}, so their design base is long. Naval nuclear
propulsion nonetheless remains a demanding capability to acquire, and
only six countries operate nuclear-powered submarines~\cite{wna_ships}.
Table~\ref{tab:intro-submarines} gives the size of each national fleet at
the end of 2019, about 150 submarines in total~\cite{wna_ships}.

\begin{table}[!htbp]
  \centering
  \caption{Nuclear-powered submarines in service by country, end of 2019.}
  \label{tab:intro-submarines}
  \small
  \begin{tabular}{@{}lr@{}}
    \toprule
    Country & Submarines \\
    \midrule
    United States  & 70 \\
    Russia         & 40 \\
    China          & 19 \\
    United Kingdom & 10 \\
    France         & 9  \\
    India          & 3  \\
    \bottomrule
  \end{tabular}

  \vspace{2pt}
  {\footnotesize Data from~\cite{wna_ships}.}
\end{table}
\FloatBarrier

Before a naval reactor goes to sea, it is common practice to prove it in
a prototype. Work on nuclear marine propulsion began in the 1940s, and
the first test reactor built for it started up in the United States in
1953. India operated a 20~MW prototype unit from 2003 for its
submarine programme~\cite{wna_ships}. The United States programme
still requires a new fuel to run in a prototype reactor before it is
deployed in the fleet~\cite{jason_2016}. Chapter~\ref{ch:litreview}
documents the naval lineages in more detail, and the Brazilian programme
follows the same path with a land-based prototype.

In Brazil, the naval nuclear propulsion programme rests on LABGENE
(\emph{Laboratório de Geração de Energia Nucleoelétrica}), a
land-based prototype plant under construction at the Centro Industrial
Nuclear de Aramar in Iperó, São Paulo~\cite{ctmsp_visita_2018}. It
was built to test the nuclear reactor of the submarine Álvaro
Alberto~\cite{vasconcelos_fapesp_2025}.

The plant serves to validate the design conditions and to test the
operating conditions of this nuclear propulsion
plant~\cite{ctmsp_labgene_2015}. The land-based prototype reactor, the
submarine, its reactor and its fuel are designed and built in
Brazil~\cite{brazil_npt_wp13_2026}.

The core data of LABGENE are closed. The open record fixes four facts
about the reactor:

\begin{enumerate}
  \item \textbf{The reactor type:} It is a pressurised water
        reactor~\cite{ctmsp_labgene_2015}.

  \item \textbf{The rated power:} According to the Navy the reactor
        generates 48\,MW of thermal power~\cite{demartini_nuclep_2022}.
        The academic literature gives 48 to 50\,MW thermal for the
        submarine reactor~\cite{dinizcosta_2017}.

  \item \textbf{The fuel class:} The fuel is low-enriched uranium (LEU),
        as declared by Brazil at the 2026 Review Conference of the
        Non-Proliferation Treaty~\cite{brazil_npt_wp13_2026}, that is,
        uranium enriched below 20\,wt\% in
        $^{235}$U~\cite{nrc_10cfr110_2}. The enrichment reported in open
        sources ranges from 6 to 8\,wt\%~\cite{dinizcosta_2017} to
        19\,wt\%~\cite{vasconcelos_fapesp_2025}.

  \item \textbf{The vessel envelope:} A structural dissertation on the
        submarine reactor, proposed by CTMSP under a research agreement
        with the University of São Paulo, gives the vessel dimensions
        adopted according to CTMSP. The cylindrical inner radius is
        900\,mm~\cite{maruyama_2013}.
\end{enumerate}

Two further points complete what is known:

\begin{itemize}
  \item A refuelling interval of five years is reported~\cite{dinizcosta_2017}.

  \item No lattice pitch, assembly count, enrichment zoning, burnable
        absorber loading, soluble boron concentration, or peaking factor
        is available from any open source.
\end{itemize}

The study is therefore built exclusively from open sources, and the
four facts above anchor the model. The other inputs are obtained in two
ways:

\begin{itemize}
  \item The fuel rod and assembly geometry are reconstructed from the
        open literature on comparable reactors, in particular the
        NuScale-class fuel design~\cite{nuscale_htp2_2017}.

  \item Every remaining parameter is a declared assumption of this
        thesis, documented with its basis in
        Appendix~\ref{app:model-inputs}.
\end{itemize}

The core studied is for that reason a LABGENE-class core, with no claim
of correspondence to the prototype itself. The thesis also tests how
far open-source tools and open data can carry a realistic naval core
design study.

The absence of a public specification also decides the form of the
study. With four inputs fixed and everything else open, the useful
question is which core configurations satisfy the known requirements.
The study is therefore an optimisation over a declared design space,
in place of the reproduction of a known core.

\section{Problem statement}
\label{sec:intro-problem}

A long refuelling interval and a flat power distribution pull a PWR
core design in opposite directions, for the following reasons:

\begin{enumerate}
  \item A multi-year cycle under the LEU ceiling needs high enrichment,
        and therefore a high beginning-of-life (BOL) reactivity.

  \item That reactivity must be held down by burnable absorbers, by
        soluble boron and by the control rods.

  \item The enrichment zoning and the absorber placement that the long
        cycle motivates degrade the radial peaking factor
        $F_{\Delta H}$, the thermal-margin quantity of first interest
        in this work.
\end{enumerate}

The design problem therefore has two competing objectives, and no
single design is best in both. Two different calculations take part in
solving it, and their costs
are very different:

\begin{itemize}
  \item \textbf{The evaluation of a design:} The OpenMC
        code~\cite{romano_openmc_2015} runs a coupled neutron transport
        and depletion calculation on the core model. One evaluation
        takes fifteen to twenty-five minutes. Because OpenMC is a Monte
        Carlo code, its results also carry statistical noise.

  \item \textbf{The search for the best designs:} A genetic algorithm
        such as NSGA-II improves a population of designs over many
        generations. In one run it asks for the objectives of several
        thousand designs.
\end{itemize}

If every request of the genetic algorithm were answered by an OpenMC
evaluation, a single search would need several thousand runs of fifteen
to twenty-five minutes each. For the 4800 evaluations of one search
described in Section~\ref{sec:bg-moo}, that is 1200 to 2000 hours, or
seven to twelve weeks of continuous computation, and the loop repeats
the search at every iteration. Coupling the search directly to OpenMC is
therefore out of reach.

In this work, the genetic algorithm never calls OpenMC. It searches a
Gaussian-process surrogate, a statistical model fitted to the designs
already evaluated with OpenMC, whose predictions cost almost nothing.
OpenMC is run only on the small number of designs that the search
proposes. Their results join the designs already evaluated, and the
Gaussian-process model is fitted again on the enlarged set, so every
pass of this active-learning loop gives an improved surrogate. The loop
continues until its stopping criterion is met or the user of the
process decides to interrupt it (Section~\ref{sec:bg-al}).

The question the thesis answers is how such a surrogate-assisted
optimisation can be built so that its Pareto front can be trusted, when
every OpenMC evaluation costs that much and carries statistical noise.

\section{Objectives}
\label{sec:intro-objectives}

The work has two objectives. The first is methodological. The second
is the application that motivated it, and it cannot be reached without
the first.

The first objective is to build a surrogate-assisted framework that can
optimise a nuclear core. The framework has three parts:

\begin{itemize}
  \item A parametric OpenMC model of the core, built at pin, assembly
        and 32-assembly fidelity, which evaluates the designs.

  \item A Gaussian-process surrogate~\cite{rasmussen_gp_2006}, which
        learns from the OpenMC evaluations.

  \item The Non-dominated Sorting Genetic Algorithm II
        (NSGA-II)~\cite{deb_nsga2_2002}, which searches the surrogate.
\end{itemize}

The three parts are joined in an active-learning loop. At each
iteration, NSGA-II searches the Gaussian-process surrogate, and from the
candidate designs it returns a small batch, six designs in this work,
is selected where the surrogate is most uncertain and run in OpenMC
(Section~\ref{sec:bg-al}). A search that would need thousands of
evaluations therefore costs only the OpenMC runs of the selected
designs, at about twenty minutes each.

Before its results can be used, the framework has to be validated in
three ways:

\begin{enumerate}
  \item The surrogate against the OpenMC evaluations it was trained on.

  \item The cheap assembly-level proxy against the core-level quantity
        it stands for.

  \item The candidate designs against a model of higher fidelity than
        the one used in the loop.
\end{enumerate}

The second objective is to obtain, from the little open data available,
a set of candidate core configurations for LABGENE, from which a final
design can be chosen once their confirmation is complete. The optimisation maximises the
cycle length, in effective full-power days (EFPD), and minimises the
radial peaking factor. Its constraints are motivated by licensing
practice:

\begin{itemize}
  \item The core must fit inside the vessel.

  \item The excess reactivity at beginning of life must not exceed the
        reactivity worth available from the burnable absorber and the
        control rods.

  \item Its enrichment must stay below the LEU ceiling.
\end{itemize}

The cycle length measures how long the core sustains full power before
it can no longer be kept critical. The peaking factor compares the most
heavily loaded fuel rod with the average rod, so a value close to unity
indicates a flat radial power distribution. A flatter distribution keeps
the hottest rod further from its thermal limits, above all the departure
from nucleate boiling, and the margin gained can be spent on power
output or on cycle length (Section~\ref{sec:bg-peaking}). Both
quantities are defined formally in Chapter~\ref{ch:background}, and
the estimators actually used are fixed in Chapter~\ref{ch:methodology}.

Because the two objectives conflict, the result of the optimisation is a
Pareto front: a set of designs in which none is better than another in
both objectives. The final design is then chosen from this front using
criteria outside the optimisation, such as controllability and the
results of higher-fidelity checks.

Campaigns 1 to 8 pose the problem in that form, and their results
changed it. Every design on the Campaign 8 front already exceeds a
refuelling interval of five years. Pushing the cycle length further
only raised the excess reactivity at beginning of life, which the
control rods then had to hold with their banks inserted deep into the
core (Section~\ref{sec:concl-reform}). The design is therefore limited
by the control of that excess reactivity, and the cycle length only has
to meet the five-year requirement. Campaign 9 poses the problem again
with two changes:

\begin{itemize}
  \item The cycle length becomes a constraint with a minimum of 1826\,EFPD,
five years at full capacity.

  \item The soluble boron concentration that the core needs at its most
        reactive state becomes the second objective.
\end{itemize}

Both objectives were reached. The framework ran nine optimisation campaigns, which together performed
636 evaluations of the Monte Carlo model, summed over all the campaigns.
Each campaign changed one element of the method relative to its
predecessor, and the framework was validated in the three ways listed
above.

The two final campaigns gave the following results:

\begin{itemize}
  \item Campaign 8 produced a front of eight feasible designs. Its
        candidate, C8-47, passed every check that was run on it.
        The ordering of its nearest neighbours on the peaking axis is
        inside the seed noise.

  \item Campaign 9 produced a set of five candidate designs on its
        front, listed in order of increasing boron demand:
        C9-47, C9-44, C9-34, C9-40 and C9-35. All five meet the
        five-year requirement with 1406 to 1908\,ppm of soluble boron
        at their most reactive state. Re-solved on the
        three-dimensional model with all rods out, C9-47, C9-34 and
        C9-40 stay on the front. C9-44 leaves it, and C9-35 drops out on
        a peaking tie that lies within the statistical uncertainty. 
        Every front member clears the boron ceiling of its own
        lattice and is held subcritical by the first two regulating
        banks on the three-dimensional model. A few smaller checks were
        still open when this thesis was written.
\end{itemize}

\section{Contributions}
\label{sec:intro-contributions}

The thesis makes two contributions. The first is a pair of Pareto fronts of possible designs for a
LABGENE-class 48\,MWth core. The core has a fixed $17\times17$ lattice,
and the design variables are the following:

\begin{itemize}
  \item One assembly enrichment, set so that no ring of the zoned core
        exceeds the low-enriched-uranium limit of 19.75\,wt\%.

  \item The gadolinia concentration.

  \item The number of poisoned pins.

  \item The reflector thickness, bounded by the vessel.
\end{itemize}

The fronts give the following results:

\begin{itemize}
  \item Every design on the front of Campaign 8
        (Section~\ref{sec:res-c8-front}) fits inside the open vessel
        envelope, holds its beginning-of-life reactivity within the
        measured worth of the regulating banks, and stays below the LEU
        ceiling.

  \item A second set from the same campaign is controllable with two
        banks alone, as a commercial PWR is operated
        (Section~\ref{sec:res-c8-split}). The comparison shows that the
        front exceeds the five-year requirement at the cost of excess
        reactivity held by deeply inserted banks, so the cycle length is
        a minimum requirement and the excess reactivity is the quantity
        to optimise (Section~\ref{sec:concl-reform}).

  \item The front of Campaign 9 (Section~\ref{sec:res-c9-front})
        answers the problem posed that way, with five designs that meet
        the five-year requirement.
\end{itemize}

The second contribution is a process that can be used in a design
project. It rests on the following rules:

\begin{enumerate}
  \item The model can be built from the basic information of a design
        project. Each input takes its value from a source where one
        exists, or from a declared assumption where it does not, and
        every input is recorded with its basis in a provenance record
        kept apart from the results.

  \item The same set of design parameters defines the model at pin,
        assembly and full-core level, so a design screened on the single
        assembly is exactly the design later confirmed on the core.

  \item The surrogate-assisted loop stops on a rule based on the
        hypervolume gain.

  \item The campaigns change one element at a time, so every difference
        in outcome between two campaigns is attributable to one cause.

  \item Constraints are measured on the model, as the controllability
        screen is, in place of bounds derived from assumptions, and the
        critical boron concentration of every design can be measured
        inside the evaluator. The results of each campaign are compared
        with the design rules and the controllability requirements, and
        that comparison feeds back into the next campaign. It either
        corrects an earlier assumption or shows that the formulation of
        the problem must change for a requirement to be met
        (Chapter~\ref{ch:methodology}).

  \item A proxy, a cheaper quantity computed in place of the one of real
        interest, such as the peaking factor of a single assembly in
        place of that of the whole core, is used only after its
        agreement with the quantity it replaces has been measured, and it is checked again
        whenever results of higher fidelity become available. A proxy
        that no longer orders the designs correctly is replaced by the quantity it represents.

  \item The candidates are confirmed with replicated random seeds and on
        a three-dimensional core model that includes the axial structure
        of the core.
\end{enumerate}

The process is documented in Chapter~\ref{ch:methodology} and
Appendix~\ref{app:code-architecture}. The published surrogate-assisted core studies reviewed in
Chapter~\ref{ch:litreview} use neural network surrogates on loading
pattern problems for cores of fixed geometry, trained before the
search. The process here searches the fuel and the geometry of a core
without a public specification. It also measures the fidelity of its
own proxies before relying on them.

\section{Scope and limitations}
\label{sec:intro-scope}

The optimisation loop runs on a two-dimensional model, with the
following treatment of the core:

\begin{itemize}
  \item Depletion is computed on the single assembly under reflective
        boundaries and is corrected for leakage through a calibrated
        reactivity target.

  \item From Campaign 4 onward, the beginning-of-life peaking objective
        is evaluated on the full 32-assembly core.

  \item From Campaign 8, the axial leakage enters the target as a
        constant measured on a finite-height model.
\end{itemize}

Three-dimensional calculations are limited. For the Campaign 8 candidates,
the three-dimensional core model gives the peaking, the rod worth, and the
axial leakage with the rods withdrawn and inserted. These calculations
showed that the axial leakage correction changes with the number of rods
inserted and with the position of the withdrawn rods, so the constant
measured on the unrodded core does not hold for the rodded core. For the
Campaign 9 front, the same model gives the peaking, the axial power shape,
the bank worth and the axial rod worth curves. No depletion is run in
three dimensions.

The boundaries of the model are stated in Chapter~\ref{ch:methodology}
where each is adopted and collected in Chapter~\ref{ch:conclusions}.
They are the following:

\begin{enumerate}
  \item The core is a fresh, uniform, single-batch loading. Reload
        strategy, loading-pattern freedom, and multi-batch operation are
        outside the design space.

  \item Thermal-hydraulic feedback is absent. Temperatures are fixed at
        representative hot full power values, so no departure from
        nucleate boiling ratio, no fuel temperature limit and no
        transient is evaluated. The moderator temperature coefficient
        is evaluated outside the optimisation loop, as a soluble boron
        ceiling on selected designs, and is never imposed as a
        constraint (Section~\ref{sec:meth-mtc}).

  \item The depletion runs at a constant soluble boron concentration of
        1000\,ppm, with no critical boron search along the cycle and no
        boron burnout. Every reported cycle length is therefore a lower
        bound for a core operated with boron letdown. Campaign 9
        measures the critical concentration at beginning of life and at
        the most reactive state of the cycle.

  \item The controllability screen tests whether the regulating banks
        hold the fresh core subcritical at power with the boron present
        and no xenon. It is a hold-down test, and the shutdown margin
        in the cold state with the highest-worth rod withdrawn is not
        demonstrated.

  \item The statistical error of the peaking objective was bounded but
        not measured inside the loop. Measured afterwards on the
        candidates, it does not resolve the ordering of neighbouring
        front designs along that axis. The work delivers a Pareto front
        of candidates that are then evaluated at higher fidelity, so an
        uncertain ordering between neighbouring designs is manageable at
        this stage.
\end{enumerate}

\section{Thesis outline}
\label{sec:intro-outline}

The thesis is organised as follows:

\begin{itemize}
  \item Chapter~\ref{ch:litreview} reviews the naval-heritage and
        commercial small modular PWR cores documented in the open
        literature. It also reviews the development of core
        optimisation methods, from mathematical programming to
        surrogate-assisted search, and the gadolinia burnable absorber
        literature.

  \item Chapter~\ref{ch:background} assembles the theory the
        methodology rests on. It holds the formal definitions of the
        multiplication factor, the cycle length and the peaking factor,
        the Monte Carlo statistics of a maximum-based estimator, the
        rank correlation used to validate a proxy, and the components
        of the surrogate-assisted loop.

  \item Chapter~\ref{ch:methodology} presents the reactor model, the
        optimisation problem, the leakage-corrected reactivity target,
        and the active-learning loop with its validation. It then
        covers the treatment of the discharge-burnup limit, the
        controllability screen, the soluble boron with the moderator
        temperature coefficient and the boron ceiling, and the campaign
        structure, including the reformulated problem of Campaign 9.

  \item Chapter~\ref{ch:results} reports the results of the nine
        campaigns. Campaigns 1 to 7 are condensed to their main results,
        including the core-level validation that changed the course of
        the work. Campaigns 8 and 9 are reported in detail: the front and
        the two-bank subset of Campaign 8 with the post-analysis of its
        candidates, the retrospective scoring of the reformulated
        problem, and the front of Campaign 9 with its confirmation.

  \item Chapter~\ref{ch:conclusions} states what the study found and
        the reformulation of the design problem that follows from it. It
        closes with the limitations of the work and the possible
        developments that its results support.

  \item Appendix~\ref{app:model-inputs} records the provenance of every
        model input.

  \item Appendix~\ref{app:code-architecture} gives an overview of the
        simulation code and of the archive of the campaigns.
\end{itemize}
